\section{Discussion}
\label{sec:discussion}

\subsection{Principal findings}
The preregistered transport-failure hypothesis was not confirmed once a
verified estimator defect was corrected. The original analysis counted every
accepted study whose report mentioned none of the five targets as correctly
classified. Such studies were about half of the source cohort and a fifth of
the external one, so the originally reported transfer gaps of about 0.04--0.05 for the
multimodal models measured the difference in label availability between the two
report corpora rather than a change in model error. Over studies with at least
one labelled target, the multimodal transfer gaps were close to zero. Their
intervals excluded a material increase, conditional on the realised calibration
sample and training run (\cref{sec:res-uncertainty}).

That near-zero estimate is itself fragile. With predictions and frozen
thresholds held fixed, plausible choices about uncertain and unmentioned
labels, the report section used for labels, aggregation across pathologies, the
training seed and the view stratum each moved the cross-site contrast by more
than its interval width, and several changed its sign. The originally reported
findings-label ``reversal'' was produced mainly by the same defect acting on an
asymmetric cohort definition. Where both report sections label the same cell
they agree in more than 95\% of cases. Under the primary label handling, the
cross-site contrast is therefore not identified well enough to support a claim
in either direction.

Three results were more stable. Frozen per-pathology thresholds transferred a
decision rule but not an operating point: at the external site the feature-fusion
model called cardiomegaly, effusion and edema far more often, and its
specificity fell, under both label handlings examined. Realised coverage stayed
close to target with context intact, under both training seeds. Misaligned
context was worse than absent context for the late-fusion model at the source
under every specification and both seeds. Externally it held for the first seed
only.

\subsection{What the endpoint measured}
The selective-risk endpoint is an error rate over the cells a report labeller
chose to label, in the studies a policy accepted. Three distinct mechanisms can
change it between hospitals, and the data separate them only partly.
\emph{Case mix} differs: with unmentioned counted as negative, effusion and
edema are roughly twice as prevalent externally. \emph{Documentation practice}
differs. Source impressions omit all five targets in about half of studies, and
external impressions carry many more uncertain mentions. The findings section
records explicit negatives at the source but is absent from most external
studies. \emph{Labeller error} may differ too. The two sites were labelled by
CheXbert, but the external labels were produced by the dataset's publishers
with undocumented settings, and agreement between report labels and image
findings is imperfect~\cite{jain2021visualchexbert,rafferty2025limitations}.
Similar aggregate positive rates (\PosSrcImpression{} and \PosExtImpression{})
did not establish comparable labels, and different prevalence alone would not
have shown different disease definitions. What the controlled analysis does
establish is that selective missingness, meaning which cells and studies carry
a label at all, was the dominant driver of the originally reported cross-site contrast
and remains a first-order driver after correction.

\subsection{Implications for cross-site benchmarking}
The defect is easy to make and invisible in aggregate output. Any per-study
loss averaged over observed labels needs a convention for studies with none,
and a guard against division by zero silently chose one here. We suggest
the following checks for any cross-site evaluation built on report-derived
labels. Report, per site, the fraction of studies and of accepted studies
with at least one labelled target. Define and test the estimator's treatment
of unlabelled studies. Fix uncertain- and unmentioned-label handling in
advance and report the alternatives. Report pathology-specific operating
characteristics, not only aggregate error. Verify the unit at which a context
intervention is applied, and replay it to confirm its constraints. Quantify
calibration-sample uncertainty alongside the evaluation bootstrap. Where a
cross-site sign matters, obtain a reference standard that is dense and
comparable at both sites.

\subsection{Implications for deployment}
On this evidence a frozen threshold should not be assumed to keep its
operating point at a new hospital, even when its coverage holds. We did not
evaluate target-site recalibration: the protocol prohibits external tuning,
and recalibration answers a different question, whether a method can be
re-tuned rather than whether a shipped policy transports. We therefore make
no claim about whether recalibration would hide or restore any particular
failure. A decision-aware confidence score ranked errors better than the frozen
margin but held its coverage less well across sites. Better ranking at the
source did not make the operating point more transportable.

\subsection{Limitations}
\label{sec:limitations}
\emph{Labels.} All labels are automatic and report-derived. No expert
reference standard was available, and the endpoint is not a clinical
five-disease error rate. The external labels' generation settings are
undocumented. The decisive limitation for cross-site claims is that neither
label source is dense and comparable at both sites. We provide a blinded
adjudication protocol (supplement) as the route to resolving it; no
adjudication has been performed.

\emph{Post hoc analyses.} The estimator correction and every analysis labelled
R1 were done after the external results had been inspected. The correction
fixes a verified defect, but it reverses the original conclusion, and readers
should weigh it with that history in mind. Both estimators are reported. The
H3 operationalisation, the S1 implementation choices and the second training
seed were also decided after inspection.

\emph{Uncertainty.} Reported intervals condition on one calibration sample,
one C2 realisation and one training run per model. Two training seeds were
examined for every model. Calibration-sample
variability was \CalRatioMin{}--\CalRatioMax{} times the evaluation-bootstrap variability for the
corrected contrasts. Two seeds cannot establish that any result is
seed-independent. The second seed changed both multimodal models' corrected
transfer gaps by about 0.04, and removed M3's external context effect.

\emph{Interventions.} C2 was applied per image, and its donors depend on the
cohort being permuted. Removing context synthetically (C1) is not the same as
context that is naturally missing. The informativeness rule measures length,
not clinical content.

\emph{Scope.} One pair of hospitals in one direction. No compatible third
site with pre-diagnostic text and comparable labels was available to us.
Reverse transfer would require training on the external site, which the
protocol does not authorise. Demographic analyses are external-only and
exploratory because source demographics were not acquired. Images differ in
compression history at origin, although preprocessing was identical.

\emph{Protocol.} Two planning documents were written on 18 July 2026 before
any data access: a prose hypothesis registry and statistical analysis plan,
and the machine-readable registry that governed execution. The prose plan led
with AURC and pathology-specific reporting and required AP/PA, one-image,
concordant-study and uncertain-label sensitivities, and it was never formally
retired. The machine-readable registry required temperature-scaled and
split-conformal baselines and listed AUGRC as a primary threshold-free metric.
None of these was run before the results were known. Several are reported here
post hoc; the one-image-per-study and concordant-study sensitivities remain
unrun, because they need image-level predictions that were not cached.

\subsection{Future work}
The most informative next step is expert adjudication of a stratified sample
of studies at both sites, against which the label-handling alternatives can be
calibrated. A third institution with pre-diagnostic text would turn one
difference into a distribution. An informativeness rule based on clinical
content rather than length would make the context interaction interpretable.
Further training seeds for all four models would bound training variability.

\section{Conclusion}
A selective policy frozen at one hospital kept its coverage at another, but
its per-pathology operating points moved. Whether its selective error rose
could not be determined from report-derived labels whose availability differs
between the hospitals. The original conclusion that it rose was an artifact
of counting label-less studies as correct. A compound-shift protocol can
isolate context effects within a site, but cross-site risk claims require a
label endpoint that is comparable across sites. That requirement should be
established before such claims are made.
